Reported benefits of advanced residential cooling control often rest on weak baselines,
unreproducible models or single runs. This paper benchmarks five controllers on one
released model of a 95\,m$^2$ top-floor apartment in Pune, India, over a
typical year at 15-min resolution: a split-unit thermostat, SIMC-tuned PID, economic
model predictive control (E-MPC), a shielded double deep Q-network (DQN), and a hybrid
E-MPC/PI controller with a safety override. The plant couples an exactly discretized
two-node RC network with an inverter VRF model, indirect evaporative cooling, an
energy-recovery outdoor-air system, and moisture, CO$_2$ and PM$_{2.5}$ balances. All
controllers face plant--model mismatch, sensor noise and forecast error, and results are
replicated over five seeds. The thermostat was cheapest but overcooled, accumulating
\KhTh\,K\,h of occupied discomfort. Among the modulating controllers the hybrid was best: it was \RelHybPidCost\,\% cheaper than PID with \RelHybPidKh\,\% less
discomfort ($p<0.001$), and \RelHybMpcCost\,\% cheaper than E-MPC with \RelHybMpcKh\,\%
less discomfort. Ablations attribute the cost saving to the PI layer and most of the
comfort gain to the safety override, and show that E-MPC's comfort loss in a heatwave comes from its
narrow comfort margin, not from prediction error. The hybrid remained the most comfortable controller across three alternative
tariffs and a 40-sample variation of eleven assumed inputs, and beat PID on
both cost and discomfort in 95\,\% of those samples; the size of its saving varied
with the building. Indoor CO$_2$ and PM$_{2.5}$ depended on ventilation and
filtration design and were nearly identical across controllers.
